\section{Formal verification}\label{app:formal}

Theorems~\ref{thm:barrier}--\ref{thm:spatial} and Corollary~\ref{cor:half} are formally verified in Lean~4
\cite{Lean4} (version 4.35.0) with the Mathlib library \cite{Mathlib}, in the repository \cite{KAMrepo} at
commit \texttt{a994efdf}. The formal proof of Theorem~\ref{thm:resonant} is in preparation in the same
repository. This appendix records what the formal proofs establish, what they do not, and where the statements
are.

\subsection{What is verified and what is not}

The Lean statements are written in the model of Section~\ref{sec:results}: the exact Hamiltonian \eqref{eq:H}
in heliocentric positions and scaled canonical momenta, the Kepler energies, axes, actions and eccentricity
vectors computed from $(q,p)$, the AMD $C=\sum_i\Lambda_i-\abs L$, and the thresholds \eqref{eq:thresholds}
defined as minima over the action chamber. The conclusions include the existence and uniqueness of the solution
of Hamilton's equations on $[-T,T]$, and the fact that its barycentric reconstruction solves Newton's equations
with masses $(1,\mu\nu_1,\dots,\mu\nu_n)$. All constants are quantified before $\mu$ and the initial state,
exactly as in the paper. The general theorems do not assume $\max_i\nu_i=1$; the normalized forms are
corollaries.

Lean checks that each stated conclusion follows from the stated hypotheses and the definitions, using Mathlib
and the axioms listed below. It does not check:
\begin{itemize}
\item that the Lean definitions express the intended physical and mathematical notions (heliocentric
  canonical coordinates, the Kepler elements, the AMD, the thresholds, the osculating elements). This was
  argued by comparing every definition and statement with the mathematical text, in reviews by readers other
  than the writers of the formal proofs;
\item the identification of $B_{ij}$ with the critical AMD of Laskar and Petit (Lemma~\ref{lem:threshold}),
  which is not formalized;
\item the proofs of Lemmas~\ref{lem:fourier}--\ref{lem:scaling}: the formal proofs of
  Theorems~\ref{thm:planar} and~\ref{thm:spatial} use instead the finite-cutoff normal form already available in
  the library (Remark~\ref{rem:engine});
\item the literature citations and the prose of this paper.
\end{itemize}
Proposition~\ref{thm:partialnek}, the input of Theorem~\ref{thm:resonant}, is itself formally proved in the same
library; its statement was fixed in advance, with \cite[Theorem~4.4]{Nek77} as its source. The file
containing that proof, at commit \texttt{a994efdf}, contains no \texttt{sorry}; the formal proof of
Theorem~\ref{thm:resonant} depends on it.

\subsection{Statements}

Table~\ref{tab:lean-main} lists the formal statements of the main results. Names in the planar files are in the
namespace \texttt{KAM.Planetary.StrongAMD}, names in the spatial files in \texttt{KAM.Planetary.Spatial}, and
the name of Proposition~\ref{thm:partialnek} in \texttt{KAM.Nekhoroshev}. Paths are relative to
\texttt{KAM/Planetary/} except for the last row.

\begin{table}[ht]
\centering\small
\begin{tabular}{@{}l>{\raggedright\arraybackslash}p{13.4cm}@{}}
\toprule
Result & Lean declaration (file and line)\\
\midrule
Thm.~\ref{thm:barrier} & \texttt{lp0\_planar\_barrier} (\texttt{StrongAMD/PlanarSafety/Main.lean:255})\\
\quad spatial & \texttt{spatialStrongBarrier\_physical}, \texttt{exists\_spatialBarrierCompact}
  (\texttt{Spatial/Geometry.lean:83, 252})\\
Thm.~\ref{thm:planar} & \texttt{lp1\_planar\_strongAMD\_safety}, \texttt{lp1\_planar\_strongAMD\_safety\_normalized}
  (\texttt{StrongAMD/PlanarSafety/Main.lean:44, 137})\\
Thm.~\ref{thm:spatial} & \texttt{lp2\_spatial\_strongAMD\_safety}, \texttt{lp2\_spatial\_strongAMD\_safety\_normalized}
  (\texttt{Spatial/SpatialSafety/Main.lean:43, 153})\\
Cor.~\ref{cor:half} & \texttt{lp1\_planar\_strongAMD\_safety\_halfInteger}
  (\texttt{StrongAMD/PlanarSafety/Main.lean:190}), \texttt{lp2\_spatial\_strongAMD\_safety\_halfInteger}
  (\texttt{Spatial/SpatialSafety/Main.lean:210})\\
Thm.~\ref{thm:resonant} & \texttt{lp3\_planar\_strongAMD\_safety}, \texttt{lp3\_spatial\_strongAMD\_safety}, and
  their normalized forms (in preparation)\\
Prop.~\ref{thm:partialnek} & \texttt{partialConvex\_nekhoroshev}
  (\texttt{KAM/Frozen/Nekhoroshev/Theorems/PartialConvexNekhoroshev.lean:48})\\
\bottomrule
\end{tabular}
\caption{Formal statements of the main results, at commit \texttt{a994efdf}.}\label{tab:lean-main}
\end{table}

Several lemmas have formal counterparts that can be read independently; Table~\ref{tab:lean-lemmas} lists the
principal ones. They formalize the content of the lemma, sometimes split differently.

\begin{table}[ht]
\centering\small
\begin{tabular}{@{}l>{\raggedright\arraybackslash}p{14.4cm}@{}}
\toprule
Lemma & Lean declarations (file and line)\\
\midrule
\ref{lem:threshold} & \texttt{exists\_pairThreshold\_attained} (\texttt{StrongAMD/Threshold/Pair.lean:97}),
  \texttt{continuousOn\_strongThreshold} (\texttt{StrongAMD/Threshold/Strong.lean:120})\\
\ref{lem:moving} & \texttt{strongAMD\_le\_of\_tube}, \texttt{maximal\_global\_of\_barrier}
  (\texttt{StrongAMD/Continuation.lean:43, 72})\\
\ref{lem:chart} & \texttt{KAM.Planetary.poincare\_symplectic} (\texttt{Poincare/Symplectic.lean:642}),
  \texttt{KAM.Planetary.poincare\_image} (\texttt{Poincare/Bijective.lean:320}),
  \texttt{lpChart\_strongAMD} (\texttt{StrongAMD/PlanarDictionary.lean:156})\\
\ref{lem:domain} & \texttt{lemma8\_complete} (\texttt{StrongAMD/AnalyticDomain/Bounds.lean:185})\\
\ref{lem:exit} & \texttt{strongFirstExit\_of\_stopped\_drift} (\texttt{StrongAMD/Confinement.lean:170})\\
\ref{lem:osc} & \texttt{osculating\_barrier} (\texttt{StrongAMD/Osculating.lean:268})\\
\ref{lem:spchart} & \texttt{spatialPoincare\_image} (\texttt{Spatial/Chart/Image.lean:180}),
  \texttt{spatialPoincare\_symplectic} (\texttt{Spatial/Chart/Symplectic.lean:256}),
  \texttt{spatialFixedPlaneDeficit\_spatialPoincareProd} (\texttt{Spatial/Product/Deficit.lean:110})\\
\ref{lem:spdomain} & \texttt{lemma18\_complete} (\texttt{Spatial/AnalyticDomain/Bounds.lean:191})\\
\bottomrule
\end{tabular}
\caption{Formal counterparts of lemmas, at commit \texttt{a994efdf}.}\label{tab:lean-lemmas}
\end{table}

\subsection{Axioms and status}

As reported by the project's build checks at commit \texttt{a994efdf}, the declarations of
Table~\ref{tab:lean-main} for Theorems~\ref{thm:barrier}--\ref{thm:spatial} and Corollary~\ref{cor:half}
depend only on the axioms \texttt{propext}, \texttt{Classical.choice} and \texttt{Quot.sound}, the standard
axioms of Lean and Mathlib; in particular they use no \texttt{sorry}. The same is reported for the development
version of the formal proof of Theorem~\ref{thm:resonant}.

Theorems~\ref{thm:barrier}--\ref{thm:spatial} and Corollary~\ref{cor:half} are proved, and their formal
statements and proofs have been independently reviewed. Theorem~\ref{thm:resonant} is proved on a development
version; its independent review is in progress. The exact statement texts of all results are under final
review.

\subsection{Use of AI and formal tools}

% FOR OWNER: confirm wording
The proofs in this paper, their Lean formalization and a first draft of the text were developed with
substantial use of large language model systems, working under the direction of the authors. Every formal proof
was checked by the Lean kernel. Statements, proofs and the correspondence between the paper and the formal
statements were reviewed by independent model-based audits and by the authors. The authors take full
responsibility for the content of the paper.
